%% =====================================================================
%%  Unit 13 — What are you doing?  (Dire ce qui se passe)
%%  Pistes : 13.01–13.04 dialogues · 13.05–13.06 vocabulaire · 13.07–13.08 grammaire
%%           13.09 prononciation · 13.10–13.13 exercices (textes en annexe)
%% =====================================================================

\newcommand{\otIllusUnXIII}{\otIllusDuo{working/104/86/15/0.10,talking/190/92/20/0.09,waiting/122/14/13/0.08,
  calling/150/8/11/0.08,reading/196/22/12/0.07,wearing/100/50/12/0.07}{emma}{What are you doing?}{lea}{I'm coming!}}

\unite[What are\\you doing?]{13}{What are you doing?}{Dire ce qui se passe}{otIllusUnXIII}

\situation{Épisode 13}{Vendredi, 15 heures au studio. Le téléphone sonne et Léa est seule à
l'accueil : tout le monde est occupé ! Le soir, c'est le grand concert du groupe de Sam… et Léa est
coincée dans le métro.}

\begin{objectifs}
\item dire ce que vous faites en ce moment et ce que font les autres ;
\item poser des questions avec \en{What are you doing?} et y répondre ;
\item distinguer le présent simple (les habitudes) et le présent en \en{-ing} (maintenant) ;
\item répondre au téléphone, demander quelqu'un et prendre un message ;
\item prononcer correctement le son \phon{ŋ} de \en{-ing}.
\end{objectifs}

\auprogramme
  {Décrire une scène, parler au téléphone}
  {Le présent en \en{be + -ing}\par\vspace{0.6mm}Présent simple ou présent en \en{-ing} ?}
  {Actions, vêtements, téléphone : 31 expressions}
  {Le son \phon{ŋ} de \en{-ing}}
  {La politesse au téléphone}

\semaine{casque/otNavy/Dialogues/Écouter et lire les deux scènes,
  bulle/otGreenDark/Vocabulaire/Apprendre les 31 expressions,
  blocs/otViolet/Grammaire/Le présent en -ing,
  sons/otOrangeDark/Sons et culture/Le son -ing et le téléphone,
  crayon/otNavy/Exercices/12 exercices corrigés,
  coche/otGreenDark/Bilan/L'essentiel et le test express}

\clearpage
%% ---------------------------------------------------------------------
\rubrique{ecouter}{Allô, j'écoute !}

\begin{dialogue}{1}{Can I take a message?}{Vendredi, 15 heures, au studio. Le téléphone sonne : Léa
  décroche.}{13.01}{13.02}
\dit{lea}{Bright Ideas, good afternoon. \hl{Léa speaking}.\note{1}}{Bright Ideas, bonjour. Léa à l'appareil.}
\dit{davis}{Hello, this is Rick Davis, from Lumen Books. \hl{Can I speak to} Mr.\ Carter, please?\note{2}}{Bonjour, ici Rick Davis, de Lumen Books. Pourrais-je parler à Mr.\ Carter, s'il vous plaît ?}
\dit{lea}{I'm sorry, he's in a meeting. \hl{He's talking} to a client.\note{3}}{Je suis désolée, il est en réunion. Il parle avec un client.}
\dit{davis}{Oh. And Emma? \hl{Is she working} today?}{Ah. Et Emma ? Elle travaille aujourd'hui ?}
\dit{lea}{Yes, she is, but she's on another call. \hl{Can I take a message}?\note{4}}{Oui, mais elle est en ligne. Puis-je prendre un message ?}
\dit{davis}{Yes, please. Can Mr.\ Carter \hl{call me back}?}{Oui, s'il vous plaît. Mr.\ Carter peut-il me rappeler ?}
\dit{lea}{Of course. Can you spell your last name, please?}{Bien sûr. Pouvez-vous épeler votre nom de famille, s'il vous plaît ?}
\dit[Davis: D, A, V, I, S. My number is two one two, five five five, oh one eight nine.]{davis}{Davis: D, A, V, I, S. My number is 212-555-0189.}{Davis : D, A, V, I, S. Mon numéro est le 212 555 0189.}
\dit{lea}{Thank you, Mr.\ Davis. \hl{I'm writing} it down.}{Merci, monsieur Davis. Je le note.}
\end{dialogue}

\begin{comprendre}
\pointnote{1} \en{\textbf{Léa speaking.}} Au téléphone, « Léa à l'appareil ». Pour se présenter, on
  dit \en{This is Léa} (et jamais \en{I am Léa}).
\pointnote{2} \en{\textbf{Can I speak to…?}} La formule polie pour demander quelqu'un au téléphone :
  « Pourrais-je parler à… ? ».
\pointnote{3} \en{\textbf{He's talking to a client.}} Le présent en \en{-ing} décrit ce qui se passe
  \textbf{en ce moment} : \en{be} + verbe + \en{-ing}. Le français dit « il parle » ou « il est en
  train de parler ».
\pointnote{4} \en{\textbf{Can I take a message?}} « Puis-je prendre un message ? ». Celui qui appelle
  dit : \en{Can I leave a message?} (« Puis-je laisser un message ? »).
\end{comprendre}

\begin{dialogue}{2}{Where are you?}{Vendredi, 19 h 50. Le concert commence à 20 heures. Emma est déjà
  dans le bar ; Léa, elle, est encore dans le métro.}{13.03}{13.04}
\dit{emma}{Léa! Where are you? \hl{What are you doing}?\note{5}}{Léa ! Où es-tu ? Qu'est-ce que tu fais ?}
\dit{lea}{\hl{I'm waiting} for the train. It's very late!}{J'attends le métro. Il a beaucoup de retard !}
\dit{emma}{Hurry up! \hl{Sam is looking} for you.}{Dépêche-toi ! Sam te cherche.}
\dit{lea}{Oh no! What's he doing now?}{Oh non ! Qu'est-ce qu'il fait, là ?}
\dit{emma}{\hl{He's walking} up and down. And Ben \hl{is eating} all the chips!}{Il fait les cent pas. Et Ben mange toutes les chips !}
\dit{lea}{Ha ha! Is it busy?}{Ha ha ! Il y a du monde ?}
\dit{emma}{Yes! A lot of people \hl{are arriving}. And \hl{you're wearing} your red dress, right?}{Oui ! Plein de gens arrivent. Et tu portes ta robe rouge, hein ?}
\dit{lea}{No, \hl{I'm not}! I'm wearing jeans and sneakers!\note{6}}{Non ! Je porte un jean et des baskets !}
\dit{emma}{Oh, Léa! OK, see you soon!}{Oh, Léa ! Bon, à tout à l'heure !}
\end{dialogue}

\begin{comprendre}
\pointnote{5} \en{\textbf{What are you doing?}} Pour la question, on place \en{be} devant le sujet :
  \en{What \textbf{are} you doing?}, \en{What \textbf{is} he doing?}
\pointnote{6} \en{\textbf{No, I'm not! I'm wearing…}} La négation : \en{be} + \en{not} + verbe en
  \en{-ing}. Pour les vêtements qu'on a sur soi aujourd'hui, on emploie le présent en \en{-ing} :
  \en{I'm wearing jeans}.
\end{comprendre}

%% ---------------------------------------------------------------------
\rubrique{mots}{Les mots de l'unité}

Voici les 31 expressions de l'unité. Attention aux petits mots qui accompagnent certains verbes :
on \en{waits \textbf{for}} quelqu'un, on \en{looks \textbf{for}} ses clés, on \en{talks \textbf{to}} un ami.

\begin{vocabulaire}{Actions et vêtements}{13.05}
\categorie{Les actions}
\mot{wait for}{ˈweɪt fɚ}{attendre}
\mot{look for}{ˈlʊk fɚ}{chercher}
\mot{talk to}{ˈtɑːk tə}{parler à}
\mot{call}{kɑːl}{appeler, téléphoner à}
\mot{arrive}{əˈraɪv}{arriver}
\mot{leave}{liːv}{partir}
\mot{sit}{sɪt}{être assis, s'asseoir}
\mot{wear}{wer}{porter (un vêtement)}
\mot{carry}{ˈkæri}{porter (dans ses bras)}
\mot{hurry up}{ˌhɝːi ˈʌp}{se dépêcher}
\categorie{Vêtements et accessoires}
\mot{a suit}{ə ˈsuːt}{un costume}
\mot{a tie}{ə ˈtaɪ}{une cravate}
\mot{sneakers}{ˈsniːkɚz}{des baskets \registre{américain}}
\mot{boots}{buːts}{des bottes}
\mot{a cap}{ə ˈkæp}{une casquette}
\mot{a hoodie}{ə ˈhʊdi}{un sweat à capuche}
\end{vocabulaire}

\begin{vocabulaire}{Au téléphone}{13.06}
\categorie{Répondre et demander}
\mot{Léa speaking.}{ˈleɪə ˌspiːkɪŋ}{Léa à l'appareil.}
\mot{Who's calling, please?}{ˌhuːz ˈkɑːlɪŋ ˌpliːz}{C'est de la part de qui ?}
\mot[Can I speak to Jack?]{Can I speak to…?}{kən aɪ ˈspiːk tə}{Pourrais-je parler à… ?}
\mot{Hold on, please.}{ˌhoʊld ˈɑːn ˌpliːz}{Ne quittez pas.}
\mot{He's in a meeting.}{hiz ɪn ə ˈmiːtɪŋ}{Il est en réunion.}
\mot{She's on another call.}{ʃiz ɑːn əˌnʌðɚ ˈkɑːl}{Elle est en ligne.}
\categorie{Les messages}
\mot{Can I take a message?}{kən aɪ ˌteɪk ə ˈmesɪdʒ}{Puis-je prendre un message ?}
\mot{Can I leave a message?}{kən aɪ ˌliːv ə ˈmesɪdʒ}{Puis-je laisser un message ?}
\mot{call back}{ˌkɑːl ˈbæk}{rappeler}
\mot{a text}{ə ˈtekst}{un texto, un SMS}
\mot{voicemail}{ˈvɔɪsmeɪl}{la messagerie vocale}
\categorie{Maintenant}
\mot{now}{naʊ}{maintenant}
\mot{right now}{ˌraɪt ˈnaʊ}{en ce moment même}
\mot{at the moment}{ət ðə ˈmoʊmənt}{en ce moment}
\mot{today}{təˈdeɪ}{aujourd'hui}
\end{vocabulaire}

%% ---------------------------------------------------------------------
\rubrique[Le présent en -ing]{grammaire}{Le présent en \en{-ing}}

\otIntertitre[otViolet]{1 · \en{be} + verbe en \en{-ing}}
\begin{center}\otTableau
\begin{tabular}{@{}l@{\hspace{8mm}}>{\enOn}l<{\enOff}@{\hspace{8mm}}l@{}}
\otTH{Forme} & \otTH{Exemple} & \otTH{Français}\\\midrule
affirmation      & I'\textbf{m} work\textbf{ing}. / She'\textbf{s} work\textbf{ing}. / They'\textbf{re} work\textbf{ing}. & je travaille…\\
négation         & I'\textbf{m not} working. / He \textbf{isn't} working.      & je ne travaille pas…\\
question         & \textbf{Are} you working? / \textbf{Is} she working?        & tu travailles ?\\
réponses courtes & Yes, I \textbf{am}. / No, she \textbf{isn't}.               & oui / non\\\bottomrule
\end{tabular}\end{center}

\begin{regle}[Quand l'employer ?]
Le présent en \en{-ing} décrit une action \textbf{en cours}, au moment où l'on parle : \en{Look! Sam is
playing the guitar.} C'est le « être en train de » du français. N'oubliez jamais \en{be} : on dit
\en{I'm working}, et non \en{I working}.
\end{regle}

\blocaudio{13.07}{Écoutez et répétez}{Six phrases au présent en \en{-ing}.}

\otIntertitre[otViolet]{2 · L'orthographe du \en{-ing}}
\begin{center}\otTableau
\begin{tabular}{@{}l@{\hspace{8mm}}>{\enOn}l<{\enOff}@{}}
\otTH{Règle} & \otTH{Exemples}\\\midrule
cas général : on ajoute \en{-ing}               & work → working · play → playing · read → reading\\
\en{-e} final muet : il disparaît                & make → making · write → writing · come → coming\\
voyelle + consonne finale : on la double         & sit → sitting · run → running · swim → swimming\\\bottomrule
\end{tabular}\end{center}

\otIntertitre[otViolet]{3 · Présent simple ou présent en \en{-ing} ?}
\begin{center}\otTableau
\begin{tabular}{@{}p{78mm}@{\hspace{8mm}}p{78mm}@{}}
\otTH{Présent simple : les habitudes} & \otTH{Présent en \en{-ing} : maintenant}\\\midrule
\enb{Jack usually \textbf{wears} jeans.} & \enb{Today, he\textbf{'s wearing} a suit!}\\[0.8mm]
\enb{Léa \textbf{works} at Bright Ideas.} & \enb{Right now, she\textbf{'s talking} on the phone.}\\[0.8mm]
\enb{Sam \textbf{plays} the guitar every day.} & \enb{Listen! He\textbf{'s playing} now.}\\[1mm]
{\fontsize{9}{11}\selectfont\color{otMuted}\en{every day}, \en{usually}, \en{always}, \en{on Sundays}} &
{\fontsize{9}{11}\selectfont\color{otMuted}\en{now}, \en{right now}, \en{at the moment}, \en{today}, \en{Look!}, \en{Listen!}}\\\bottomrule
\end{tabular}\end{center}

\begin{regle}[Les verbes qui refusent le \en{-ing}]
Certains verbes décrivent un état, pas une action : \en{like}, \en{love}, \en{want}, \en{need},
\en{know}, \en{understand}. Ils restent au présent simple, même maintenant : \en{I want a coffee
now}, \en{Do you understand?}
\end{regle}

\blocaudio{13.08}{Écoutez et répétez}{Une question au présent en \en{-ing}, puis sa réponse.}

\begin{attention}[Les pièges des francophones]
\fauxjuste{I working now.}{I'm working now.}
\fauxjuste{What you are doing?}{What are you doing?}
\fauxjuste{He's siting on the sofa.}{He's sitting on the sofa.}
\fauxjuste{I'm wanting a coffee.}{I want a coffee.}
\fauxjuste{Are you understanding?}{Do you understand?}
\end{attention}

%% ---------------------------------------------------------------------
\rubrique[Le son -ing]{sons}{Le son \phon{ŋ} de \en{-ing}}

Bonne nouvelle : vous connaissez déjà ce son, grâce aux mots français \textit{parking},
\textit{camping} ou \textit{shopping}. Le son \phon{ŋ} se forme au fond de la bouche : l'arrière de la
langue touche le palais, et l'air sort par le nez. Deux erreurs à éviter : dire « inne » \phon{ɪn}, ou
ajouter un \phon{ɡ} à la fin (« sing-gue »). Dans \en{singing}, on n'entend \textbf{aucun} \phon{ɡ}.

\begin{paires}{13.09}{Écoutez chaque paire et répétez : \phon{n} (le bout de la langue) puis \phon{ŋ}
(le fond de la langue).}
\sonpaire{thi\otH{n}}{θɪn}{mince}{thi\otH{ng}}{θɪŋ}{une chose}
\sonpaire{wi\otH{n}}{wɪn}{gagner}{wi\otH{ng}}{wɪŋ}{une aile}
\sonpaire{si\otH{n}}{sɪn}{un péché}{si\otH{ng}}{sɪŋ}{chanter}
\sonpaire{ra\otH{n}}{ræn}{a couru}{ra\otH{ng}}{ræŋ}{a sonné}
\sonpaire{ba\otH{n}}{bæn}{interdire}{ba\otH{ng}}{bæŋ}{un grand bruit}
\sonpaire{su\otH{n}}{sʌn}{le soleil}{su\otH{ng}}{sʌŋ}{chanté}
\end{paires}

%% ---------------------------------------------------------------------
\rubrique{culture}{La politesse au téléphone}

\begin{culture}[Décrocher comme un pro]
À la maison, on décroche simplement par \en{Hello?}. Au travail, on annonce l'entreprise, puis son
prénom : \en{Bright Ideas, Léa speaking. How can I help you?} Pour savoir qui appelle, ne dites jamais
\en{Who are you?}, beaucoup trop brutal : préférez \en{Who's calling, please?}, ou, plus formel,
\en{May I ask who's calling?}
\end{culture}

\begin{paire}
\begin{culture}[Les textos]
Les Américains abrègent beaucoup : \en{LOL} (\en{laughing out loud}, « mdr »), \en{BRB} (\en{be right
back}, « je reviens »), \en{OMW} (\en{on my way}, « j'arrive »), \en{TTYL} (\en{talk to you later}),
\en{thx} (\en{thanks}), \en{u} (\en{you}).
\end{culture}
\begin{culture}[La messagerie]
Message type d'un répondeur : \en{Hi, this is Jack. I can't come to the phone right now. Please leave
a message after the beep!} (le bip). Beaucoup d'Américains préfèrent aujourd'hui un texto à un appel.
\end{culture}
\end{paire}

%% ---------------------------------------------------------------------
\rubrique{pratique}{Exercices}

\exercice[ecrire,niveau=1]{La forme en \en{-ing}}{Écrivez la forme en \en{-ing} de chaque verbe.}
\begin{anglais}
\begin{questions}[4]
\item work → \trou[17mm]{working}
\item make → \trou[17mm]{making}
\item sit → \trou[17mm]{sitting}
\item play → \trou[17mm]{playing}
\item write → \trou[17mm]{writing}
\item run → \trou[17mm]{running}
\item have → \trou[17mm]{having}
\item swim → \trou[17mm]{swimming}
\end{questions}
\end{anglais}

\exercice[lire,ecrire,niveau=1]{Que font-ils en ce moment ?}{Regardez chaque image et écrivez une
phrase au présent en \en{-ing}, avec \en{he} ou \en{she}.}
\begin{tcbraster}[raster columns=3,raster equal height=rows,raster column skip=3mm,raster row skip=3mm,raster before skip=1mm]
\otCarteAction{lea}{objtelephone}{otGreenDark}{talk on the phone}{She's talking on the phone.}
\otCarteAction{jack}{objtasse}{otOrange}{drink a coffee}{He's drinking a coffee.}
\otCarteAction{emma}{objlivre}{otViolet}{read a book}{She's reading a book.}
\otCarteAction{sam}{casque}{otTeal}{listen to music}{He's listening to music.}
\otCarteAction{carter}{crayon}{otOrangeDark}{draw a logo}{He's drawing a logo.}
\otCarteAction{lopez}{micro}{otRed}{sing}{She's singing.}
\end{tcbraster}

\exercice[ecrire,niveau=2]{Négations et questions}{Transformez chaque phrase : (?) en question, (−) en
phrase négative.}
\begin{questions}
\item Léa is talking on the phone. (?) \hfill \trou[70mm]{Is Léa talking on the phone?}
\item They're waiting for the bus. (−) \hfill \trou[70mm]{They aren't waiting for the bus.}
\item Sam is wearing a suit. (?) \hfill \trou[70mm]{Is Sam wearing a suit?}
\item I'm reading a book. (−) \hfill \trou[70mm]{I'm not reading a book.}
\item Emma is having lunch. (−) \hfill \trou[70mm]{Emma isn't having lunch.}
\item You're listening to me. (?) \hfill \trou[70mm]{Are you listening to me?}
\end{questions}

\exercice[ecrire,niveau=2]{Habitude ou maintenant ?}{Mettez le verbe au présent simple ou au présent
en \en{-ing}.}
\begin{questions}
\item Jack usually \trou[30mm]{wears} jeans. \hfill (\en{wear})
\item But today, he \trou[30mm]{is wearing} a suit! \hfill (\en{wear})
\item Listen! Sam \trou[30mm]{is playing} the guitar. \hfill (\en{play})
\item Sam \trou[30mm]{plays} the guitar every day. \hfill (\en{play})
\item Shh! The baby \trou[30mm]{is sleeping}. \hfill (\en{sleep})
\item I \trou[30mm]{don't drink} coffee in the evening. \hfill (\en{not drink})
\item What \trou[30mm]{are you doing} right now? \hfill (\en{you / do})
\item \trou[30mm]{Do you speak} French? \hfill (\en{you / speak})
\end{questions}

\exercice[lire,niveau=2]{Les textos de Léa}{Samedi matin, Jack envoie un texto à Léa. Lisez la
conversation, puis indiquez si chaque phrase est vraie (V) ou fausse (F).}
\begin{otTextos}{Jack}
\otSms{g}{Hi Léa! What are you doing?}
\otSms{d}{Hi Jack! I'm having breakfast. I'm so tired! What about you?}
\otSms{g}{I'm at the park with Kate and Josh. We're playing soccer. Josh is winning! LOL}
\otSms{d}{Ha ha! Is Kate playing too?}
\otSms{g}{No, she isn't. She's sitting on a bench. She's reading a book.}
\otSms{d}{Have fun! I'm going back to bed! BRB… in 2 hours!}
\end{otTextos}
\par\vspace{1mm}\noindent{\fontsize{8.6}{10}\selectfont\color{otMuted}\textit{What about you?} : et toi ? ·
\textit{winning} : en train de gagner · \textit{a bench} : un banc · \textit{Have fun!} : amuse-toi bien !\par}
\begin{anglais}
\begin{questions}[2]
\item Léa is having breakfast. \hfill\trou[9mm]{V}
\item Jack is at home. \hfill\trou[9mm]{F}
\item Jack is playing basketball. \hfill\trou[9mm]{F}
\item Josh is winning. \hfill\trou[9mm]{V}
\item Kate is playing soccer. \hfill\trou[9mm]{F}
\item Kate is reading. \hfill\trou[9mm]{V}
\end{questions}
\end{anglais}

\exercice[lire,niveau=1]{Au téléphone}{Associez chaque phrase à sa réponse : écrivez la lettre.}
\begin{anglais}
\noindent\begin{minipage}[t]{0.48\linewidth}
\begin{questions}[1][english]
\item Can I speak to Emma, please? \hfill\trou[9mm]{b}
\item Who's calling, please? \hfill\trou[9mm]{d}
\item Can I take a message? \hfill\trou[9mm]{a}
\item Is Mr.\ Carter there? \hfill\trou[9mm]{f}
\item What are you doing? \hfill\trou[9mm]{c}
\item Are you working today? \hfill\trou[9mm]{e}
\end{questions}
\end{minipage}\hfill
\begin{minipage}[t]{0.46\linewidth}
\begin{questions}[1][english]
\item[a.] Yes, please. Can she call me back?
\item[b.] Hold on, please.
\item[c.] I'm cooking dinner.
\item[d.] This is Sam Brooks.
\item[e.] No, I'm not. It's my day off.
\item[f.] I'm sorry, he's in a meeting.
\end{questions}
\end{minipage}
\end{anglais}
\par\vspace{1mm}{\fontsize{8.6}{10}\selectfont\color{otMuted}\textit{day off} : jour de congé}

\exercice[ecouter,niveau=1,piste=13.10]{\phon{n} ou \phon{ŋ} ?}{Écoutez. Pour chaque numéro, la voix
dit un seul des deux mots, deux fois. Cochez le mot entendu.}
\begin{questions}[3]
\item \choix{thin}\quad\choix{thing}
\item \choix{win}\quad\choix{wing}
\item \choix{sin}\quad\choix{sing}
\item \choix{ran}\quad\choix{rang}
\item \choix{ban}\quad\choix{bang}
\item \choix{sun}\quad\choix{sung}
\end{questions}

\exercice[ecouter,ecrire,niveau=2,piste=13.11]{Un samedi à Chicago}{Mrs.\ Lopez appelle sa fille
Sofia, à Chicago. Écoutez et complétez avec le verbe en \en{-ing}.}
\begin{anglais}
\begin{questions}
\item Sofia is \trou[30mm]{drinking} a coffee in the kitchen.
\item Leo is \trou[30mm]{reading} in his bedroom.
\item Max is \trou[30mm]{playing} in the garden with the dog.
\item David isn't \trou[30mm]{working} today.
\item He's \trou[30mm]{cooking} pancakes!
\end{questions}
\end{anglais}

\exercice[ecouter,ecrire,niveau=2,piste=13.12]{Dictée à trous}{Jack répond au téléphone du studio.
Écoutez et complétez.}
\begin{anglais}
\noindent\begin{tabular}{@{}l@{\hspace{6mm}}l@{}}
\otNomDit{jack}      & Bright Ideas, Jack \trou{speaking}.\\[1.8mm]
\otNomDit{passante}  & Hello. Can I \trou{speak} to Léa Martin, please?\\[1.8mm]
\otNomDit{jack}      & \trou{Hold} on, please\ldots{} Sorry, she's in a \trou{meeting}.\\[1.8mm]
\otNomDit{passante}  & Oh. Can I \trou{leave} a message?\\[1.8mm]
\otNomDit{jack}      & Of course. Who's \trou{calling}, please?\\[1.8mm]
\otNomDit{passante}  & \trou{This} is Anna, from the gym. Can she call me back?\\[1.8mm]
\otNomDit{jack}      & Sure. I'm \trou{writing} it down.\\
\end{tabular}
\end{anglais}

\exercice[ecrire,niveau=2]{Traduisez}{Traduisez en anglais.}
\begin{questions}[1][french]
\item Qu'est-ce que tu fais ? — Je regarde la télé.\lignes{1}
\item Elle ne travaille pas aujourd'hui.\lignes{1}
\item Écoute ! Sam joue de la guitare.\lignes{1}
\item Il porte un costume et une cravate.\lignes{1}
\item Ne quittez pas, s'il vous plaît.\lignes{1}
\item Il est en réunion. Puis-je prendre un message ?\lignes{1}
\item Nous attendons le bus.\lignes{1}
\item Tu m'écoutes ? — Oui. \textit{(réponse courte)}\lignes{1}
\end{questions}

\exercice[ecrire,niveau=3]{Un texto à un ami}{Un ami vous écrit : \en{What are you doing?} Répondez en
cinq phrases : où vous êtes, ce que vous faites, ce que font les gens autour de vous, ce que vous
portez.}
\lignes{5}

\exercice[ecouter,parler,niveau=3,piste=13.13]{À vous de parler !}{Sam vous appelle. Répondez
\textbf{à voix haute} pendant chaque pause.}
\begin{center}\otTableau
\begin{tabular}{@{}>{\otSemi\fontsize{8.6}{10}\selectfont\color{otNavy}}l@{\hspace{3mm}}>{\enOn}p{70mm}<{\enOff}@{\hspace{6mm}}p{78mm}@{}}
 & \otTH{Sam dit\ldots} & \otTH{Vous\ldots}\\\midrule
1. & Hi, it's Sam! What are you doing? & dites ce que vous faites.\\
2. & Where are you? & dites où vous êtes.\\
3. & Are you watching TV? & répondez par une réponse courte.\\
4. & What are you wearing today? & décrivez vos vêtements.\\
5. & Can you come to the café now? & acceptez ou refusez poliment.\\\bottomrule
\end{tabular}\end{center}

%% ---------------------------------------------------------------------
\rubrique{bilan}{Faites le point}

\begin{essentiel}
\begin{minipage}[t]{0.54\linewidth}
\otEssTitre{Les phrases clés}
\renewcommand{\arraystretch}{1.32}
\begin{tabular}{@{}>{\enOn\otSemi\fontsize{9.4}{11}\selectfont}l<{\enOff}@{\hspace{3mm}}>{\fontsize{9.2}{11}\selectfont\color{white!78!otNavy}}l@{}}
What are you doing?        & Qu'est-ce que tu fais ?\\
I'm waiting for the train. & J'attends le métro.\\
He isn't working today.    & Il ne travaille pas aujourd'hui.\\
Léa speaking.              & Léa à l'appareil.\\
Hold on, please.           & Ne quittez pas.\\
Can I take a message?      & Puis-je prendre un message ?\\
\end{tabular}
\end{minipage}\hfill
\begin{minipage}[t]{0.41\linewidth}
\otEssTitre{Le présent en \en{-ing}}
{\fontsize{9.2}{11.6}\selectfont \en{be} + verbe + \en{-ing} : ce qui se passe maintenant\par
\en{make} → \en{making} · \en{sit} → \en{sitting}\par}
\otEssTitre{Habitude ou maintenant ?}
{\fontsize{9.2}{11.6}\selectfont \en{He usually wears jeans.} · \en{Today, he's wearing a suit.}\par}
\otEssTitre{Prononciation}
{\fontsize{9.2}{11.6}\selectfont \en{-ing} \phon{ɪŋ}, sans \phon{ɡ} à la fin · \en{thin} ≠ \en{thing}\par}
\end{minipage}
\end{essentiel}

\begin{jesais}
\item dire ce que je fais en ce moment et ce que font les autres ;
\item poser une question au présent en \en{-ing} et y répondre ;
\item écrire correctement la forme en \en{-ing} (\en{making}, \en{sitting}) ;
\item choisir entre présent simple et présent en \en{-ing} ;
\item répondre au téléphone, demander quelqu'un et prendre un message ;
\item prononcer le son \phon{ŋ} sans ajouter de \phon{ɡ}.
\end{jesais}

\clearpage
\begin{testexpress}
\qcm[b]{« Qu'est-ce que tu fais ? » se dit :}{\en{What do you doing?}}{\en{What are you doing?}}{\en{What you are doing?}}
\qcm[a]{La forme en \en{-ing} de \en{sit} est :}{\en{sitting}}{\en{siting}}{\en{sitteing}}
\qcm[c]{« Elle ne travaille pas aujourd'hui » se dit :}{\en{She doesn't working today.}}{\en{She not working today.}}{\en{She isn't working today.}}
\qcm[a]{\en{Are you listening?} — Réponse courte affirmative :}{\en{Yes, I am.}}{\en{Yes, I'm.}}{\en{Yes, I listening.}}
\qcm[b]{« Aujourd'hui, Jack porte un costume » se dit :}{\en{Today, Jack wear a suit.}}{\en{Today, Jack is wearing a suit.}}{\en{Today, Jack is wear a suit.}}
\qcm[c]{« Je veux un café » se dit :}{\en{I'm wanting a coffee.}}{\en{I wanting a coffee.}}{\en{I want a coffee.}}
\qcm[a]{Au téléphone, « Ne quittez pas » se dit :}{\en{Hold on, please.}}{\en{Don't leave, please.}}{\en{Stay here, please.}}
\qcm[b]{Pour savoir qui appelle, on dit :}{\en{Who are you?}}{\en{Who's calling, please?}}{\en{What's your name?}}
\qcm[c]{Dans \en{singing}, on entend :}{deux \phon{ɡ}}{un \phon{n} puis un \phon{ɡ}}{\phon{ŋ}, sans \phon{ɡ}}
\qcm[b]{Dans un texto, \en{OMW} signifie :}{\en{Oh my word!}}{\en{On my way!}}{\en{One more week}}
\end{testexpress}

\noindent{\fontsize{9.6}{12}\selectfont\color{otMuted}Niveau visé (CECRL, niveau A1) : « peut décrire
simplement ce que font les gens et tenir une conversation téléphonique très simple ». Les corrigés de
l'unité sont page~\pageref{cor-u13-1} ; les textes des pistes d'exercices,
page~\pageref{trans-13.10}.\par}
